
\documentclass{article}
\usepackage{colortbl}
\usepackage{makecell}
\usepackage{multirow}
\usepackage{supertabular}

\begin{document}

\newcounter{utterance}

\centering \large Interaction Transcript for game `matchit\_ascii', experiment `same\_grid', episode 5 with NEW\_Reflect\_Agent.
\vspace{24pt}

{ \footnotesize  \setcounter{utterance}{1}
\setlength{\tabcolsep}{0pt}
\begin{supertabular}{c@{$\;$}|p{.15\linewidth}@{}p{.15\linewidth}p{.15\linewidth}p{.15\linewidth}p{.15\linewidth}p{.15\linewidth}}
    \# & $\;$A & \multicolumn{4}{c}{Game Master} & $\;\:$B\\
    \hline

    \theutterance \stepcounter{utterance}  
    & & \multicolumn{4}{p{0.6\linewidth}}{
        \cellcolor[rgb]{0.9,0.9,0.9}{
            \makecell[{{p{\linewidth}}}]{
                \texttt{\tiny{[P1$\langle$GM]}}
                \texttt{You are participating in a collaborative guessing game.} \\
\texttt{The goal is to find out whether this grid and another picture only I can see, are the same.} \\
\texttt{Please describe your grid first. Then, I will provide my description and we can ask each other questions about the grids to figure out whether they are the same. Now start your short grid description with "DESCRIPTION:" followed by the description. Do not add anything else. This is your grid:} \\
\texttt{X X ▢ X X} \\
\texttt{X X ▢ X X} \\
\texttt{▢ ▢ ▢ ▢ ▢} \\
\texttt{X X ▢ X X} \\
\texttt{X X ▢ X X} \\
            }
        }
    }
    & \\ \\

    \theutterance \stepcounter{utterance}  
    & \multicolumn{4}{p{0.45\linewidth}}{
        \cellcolor[rgb]{0.95,0.95,0.95}{
            \makecell[{{p{\linewidth}}}]{
                \texttt{\tiny{[P1$|$internal]}}
                \texttt{Reflections from previous episodes:} \\
\texttt{{-}} \\
\texttt{ Failed. I gave a clear description of my grid pattern, but the user did not provide their description or engage in questioning, ending the episode prematurely. Next time, I would follow my initial instruction exactly as given, but cannot control the other player's responsiveness. The outcome depended on collaboration that didn't occur.} \\
\texttt{{-}} \\
\texttt{ Failed. I described the grid as a symmetrical X pattern, which was correct, but I didn’t ask clarifying questions about the other player’s grid to confirm details like exact dimensions or symmetry. Next time, I would follow my description with a direct question about their grid’s structure to compare systematically.} \\
\texttt{{-}} \\
\texttt{ Failed. I gave a vague description ("grid with X's forming patterns, mostly in center columns") that didn't capture the exact shape. My partner's description didn't match, but I couldn't pinpoint differences because my initial description lacked precise row{-}by{-}row or coordinate details. Next time, I would give a structured description (e.g., row{-}by{-}row counts or positions) to enable exact comparison.} \\
            }
        }
    }
    & & \\ \\

    \theutterance \stepcounter{utterance}  
    & \multicolumn{4}{p{0.45\linewidth}}{
        \cellcolor[rgb]{0.95,0.95,0.95}{
            \makecell[{{p{\linewidth}}}]{
                \texttt{\tiny{[P1$|$internal]}}
                \texttt{{[}attempt 1 – valid{]}} \\
\texttt{DESCRIPTION: The grid is 5x5. The first, second, fourth, and fifth rows have X's in the first, second, fourth, and fifth columns, with a ▢ in the third column. The third row is entirely ▢. The pattern forms a cross shape with X's in the center column and the center row, but the center column is fully X's except for the center cell, which is ▢. Wait, no—let me correct that. Actually, the center column (third column) has X's in rows 1, 2, 4, 5 and ▢ in row 3. The center row (third row) is all ▢. So the X's form a cross with the vertical bar missing the center cell and the horizontal bar missing the center cell. The X's are in positions: (1,1), (1,2), (1,4), (1,5), (2,1), (2,2), (2,4), (2,5), (4,1), (4,2), (4,4), (4,5), (5,1), (5,2), (5,4), (5,5). The ▢s are in (1,3), (2,3), (3,1), (3,2), (3,3), (3,4), (3,5), (4,3),} \\
            }
        }
    }
    & & \\ \\

    \theutterance \stepcounter{utterance}  
    & \multicolumn{4}{p{0.6\linewidth}}{
        \cellcolor[rgb]{0.9,0.9,0.9}{
            \makecell[{{p{\linewidth}}}]{
                \texttt{\tiny{[P1$\rangle$GM]}}
                \texttt{DESCRIPTION: The grid is 5x5. The first, second, fourth, and fifth rows have X's in the first, second, fourth, and fifth columns, with a ▢ in the third column. The third row is entirely ▢. The pattern forms a cross shape with X's in the center column and the center row, but the center column is fully X's except for the center cell, which is ▢. Wait, no—let me correct that. Actually, the center column (third column) has X's in rows 1, 2, 4, 5 and ▢ in row 3. The center row (third row) is all ▢. So the X's form a cross with the vertical bar missing the center cell and the horizontal bar missing the center cell. The X's are in positions: (1,1), (1,2), (1,4), (1,5), (2,1), (2,2), (2,4), (2,5), (4,1), (4,2), (4,4), (4,5), (5,1), (5,2), (5,4), (5,5). The ▢s are in (1,3), (2,3), (3,1), (3,2), (3,3), (3,4), (3,5), (4,3),} \\
            }
        }
    }
    & & \\ \\

    \theutterance \stepcounter{utterance}  
    & & & \multicolumn{2}{p{0.3\linewidth}}{
        \cellcolor[rgb]{0.9,0.9,0.9}{
            \makecell[{{p{\linewidth}}}]{
                \texttt{\tiny{[GM$|$GM]}}
                \texttt{continue} \\
            }
        }
    }
    & & \\ \\

    \theutterance \stepcounter{utterance}  
    & & \multicolumn{4}{p{0.6\linewidth}}{
        \cellcolor[rgb]{0.9,0.9,0.9}{
            \makecell[{{p{\linewidth}}}]{
                \texttt{\tiny{[GM$\rangle$P2]}}
                \texttt{You are participating in a collaborative guessing game.} \\
\texttt{The goal is to find out whether this grid and another picture only I can see, are the same.} \\
\texttt{Please describe your grid first. Then, I will provide my description and we can ask each other questions about the grids to figure out whether they are the same. Now start your short grid description with "DESCRIPTION:" followed by the description. Do not add anything else. This is your grid:} \\
\texttt{X X ▢ X X} \\
\texttt{X X ▢ X X} \\
\texttt{▢ ▢ ▢ ▢ ▢} \\
\texttt{X X ▢ X X} \\
\texttt{X X ▢ X X} \\
            }
        }
    }
    & \\ \\

    \theutterance \stepcounter{utterance}  
    & & & \multicolumn{4}{p{0.45\linewidth}}{
        \cellcolor[rgb]{0.95,0.95,0.95}{
            \makecell[{{p{\linewidth}}}]{
                \texttt{\tiny{[P2$|$internal]}}
                \texttt{Reflections from previous episodes:} \\
\texttt{{-}} \\
\texttt{ FAIL. I did not provide a description of my grid as instructed, which prevented the comparison process from starting. Next time, I must strictly follow the initial instruction to start with "DESCRIPTION:" and give my grid's layout before any questions.} \\
\texttt{{-}} \\
\texttt{ Failed. My description was too vague ("symmetric pattern," "T shape," "U shape") and could match multiple grids. I should have given a precise structural description, like "X's are on all positions where row index equals column index or row index + column index equals 4" (the main diagonals). Next time, I will use exact mathematical or coordinate{-}based descriptions to avoid ambiguity.} \\
\texttt{{-}} \\
\texttt{ Failed. My description was inaccurate: I described a diagonal pattern, but my grid actually had a vertical "X" band in columns 2{-}4, not a true diagonal. This misled the comparison from the start. Next time, I would describe the pattern more precisely: "A 5x5 grid with a vertical block of X's in columns 2, 3, and 4, except for some missing X's forming a checkerboard{-}like absence in the center rows."} \\
            }
        }
    }
     \\ \\

    \theutterance \stepcounter{utterance}  
    & & & \multicolumn{4}{p{0.45\linewidth}}{
        \cellcolor[rgb]{0.95,0.95,0.95}{
            \makecell[{{p{\linewidth}}}]{
                \texttt{\tiny{[P2$|$internal]}}
                \texttt{{[}attempt 1 – valid{]}} \\
\texttt{DESCRIPTION: A 5x5 grid where the X's are in all positions except for column 2 (0{-}indexed), which is entirely empty (all ▢), and the rest of the grid has X's in every position except column 2. Specifically, X's are present in columns 0, 1, 3, and 4 for all rows, and column 2 is all ▢.} \\
            }
        }
    }
     \\ \\

    \theutterance \stepcounter{utterance}  
    & & & \multicolumn{4}{p{0.6\linewidth}}{
        \cellcolor[rgb]{0.9,0.9,0.9}{
            \makecell[{{p{\linewidth}}}]{
                \texttt{\tiny{[GM$\langle$P2]}}
                \texttt{DESCRIPTION: A 5x5 grid where the X's are in all positions except for column 2 (0{-}indexed), which is entirely empty (all ▢), and the rest of the grid has X's in every position except column 2. Specifically, X's are present in columns 0, 1, 3, and 4 for all rows, and column 2 is all ▢.} \\
            }
        }
    }
     \\ \\

    \theutterance \stepcounter{utterance}  
    & & & \multicolumn{2}{p{0.3\linewidth}}{
        \cellcolor[rgb]{0.9,0.9,0.9}{
            \makecell[{{p{\linewidth}}}]{
                \texttt{\tiny{[GM$|$GM]}}
                \texttt{continue} \\
            }
        }
    }
    & & \\ \\

    \theutterance \stepcounter{utterance}  
    & & \multicolumn{4}{p{0.6\linewidth}}{
        \cellcolor[rgb]{0.9,0.9,0.9}{
            \makecell[{{p{\linewidth}}}]{
                \texttt{\tiny{[GM$\rangle$P2]}}
                \texttt{This is my DESCRIPTION: The grid is 5x5. The first, second, fourth, and fifth rows have X's in the first, second, fourth, and fifth columns, with a ▢ in the third column. The third row is entirely ▢. The pattern forms a cross shape with X's in the center column and the center row, but the center column is fully X's except for the center cell, which is ▢. Wait, no—let me correct that. Actually, the center column (third column) has X's in rows 1, 2, 4, 5 and ▢ in row 3. The center row (third row) is all ▢. So the X's form a cross with the vertical bar missing the center cell and the horizontal bar missing the center cell. The X's are in positions: (1,1), (1,2), (1,4), (1,5), (2,1), (2,2), (2,4), (2,5), (4,1), (4,2), (4,4), (4,5), (5,1), (5,2), (5,4), (5,5). The ▢s are in (1,3), (2,3), (3,1), (3,2), (3,3), (3,4), (3,5), (4,3),} \\
\texttt{Now ask a question in order to find out new aspects of my grid that may be different to your grid. Start with "QUESTION:" and do not add anything else.} \\
            }
        }
    }
    & \\ \\

    \theutterance \stepcounter{utterance}  
    & & & \multicolumn{4}{p{0.6\linewidth}}{
        \cellcolor[rgb]{0.9,0.9,0.9}{
            \makecell[{{p{\linewidth}}}]{
                \texttt{\tiny{[GM$\langle$P2]}}
                \texttt{DESCRIPTION: The grid is 5x5. The first, second, fourth, and fifth rows have X's in the first, second, fourth, and fifth columns, with a ▢ in the third column. The third row is entirely ▢. The pattern forms a cross shape with X's in the center column and the center row, but the center column is fully X's except for the center cell, which is ▢. Wait, no—let me correct that. Actually, the center column (third column) has X's in rows 1, 2, 4, 5 and ▢ in row 3. The center row (third row) is all ▢. So the X's form a cross with the vertical bar missing the center cell and the horizontal bar missing the center cell. The X's are in positions: (1,1), (1,2), (1,4), (1,5), (2,1), (2,2), (2,4), (2,5), (4,1), (4,2), (4,4), (4,5), (5,1), (5,2), (5,4), (5,5). The ▢s are in (1,3), (2,3), (3,1), (3,2), (3,3), (3,4), (3,5), (4,3),} \\
            }
        }
    }
     \\ \\

    \theutterance \stepcounter{utterance}  
    & & & \multicolumn{2}{p{0.3\linewidth}}{
        \cellcolor[rgb]{0.9,0.9,0.9}{
            \makecell[{{p{\linewidth}}}]{
                \texttt{\tiny{[GM$|$GM]}}
                \texttt{abort, first word: DESCRIPTION:, but expected QUESTION:} \\
            }
        }
    }
    & & \\ \\

\end{supertabular}
}

\end{document}
